% ============================================================
%  USER GUIDE
%  Bug-Bite Classification via Tropical Cyclone Transfer Learning
%  Braude College of Engineering — Software Engineering Capstone
%
%  Usage in Overleaf:
%    \appendix
%    % ============================================================
%  USER GUIDE
%  Bug-Bite Classification via Tropical Cyclone Transfer Learning
%  Braude College of Engineering — Software Engineering Capstone
%
%  Usage in Overleaf:
%    \appendix
%    % ============================================================
%  USER GUIDE
%  Bug-Bite Classification via Tropical Cyclone Transfer Learning
%  Braude College of Engineering — Software Engineering Capstone
%
%  Usage in Overleaf:
%    \appendix
%    % ============================================================
%  USER GUIDE
%  Bug-Bite Classification via Tropical Cyclone Transfer Learning
%  Braude College of Engineering — Software Engineering Capstone
%
%  Usage in Overleaf:
%    \appendix
%    \input{user_guide}
%
%  Or compile standalone:
%    Uncomment the preamble block below, and add \end{document} at the end.
% ============================================================

% --- Standalone preamble (remove if including in main document) ----------
\documentclass[12pt,a4paper]{article}
\usepackage[utf8]{inputenc}
\usepackage[T1]{fontenc}
\usepackage{geometry}
\geometry{margin=2.5cm}
\usepackage{hyperref}
\usepackage{listings}
\usepackage{xcolor}
\usepackage{booktabs}
\usepackage{array}
\usepackage{enumitem}
\usepackage{parskip}
\lstset{basicstyle=\ttfamily\small, breaklines=true, frame=single,
        backgroundcolor=\color{gray!8}}
\begin{document}
\title{User Guide\\
       \large Bug-Bite Classification via Tropical Cyclone Transfer Learning\\
       Braude College of Engineering --- Software Engineering Capstone}
\author{Ishay Yulzary, Nathan Trostianitser}
\date{}
\maketitle
\tableofcontents
\newpage
% -------------------------------------------------------------------------

\section{User Guide}
\label{app:user-guide}

This guide covers two use scenarios for the system:

\begin{enumerate}
    \item \textbf{Running the automated training pipeline:} execute the
          full hyperparameter sweep to train and evaluate all model
          configurations. This is the primary use of the system.
    \item \textbf{Running inference via the demo notebook:} classify a
          single bug-bite image using the pre-trained ensemble on Google
          Colab, without any local setup.
\end{enumerate}

% ============================================================
\subsection{What the System Does}
\label{sec:ug-overview}
% ============================================================

The primary output of this project is an automated machine learning
pipeline that trains and evaluates an ensemble of neural network classifiers
under two paradigms:

\begin{description}
    \item[Control] Models trained directly from ImageNet weights on the
          bug-bite dataset (single-stage transfer).
    \item[TC (Transfer Classification)] Models first pre-trained on tropical
          cyclone satellite imagery, then fine-tuned on the bug-bite dataset
          (two-stage transfer). This is the primary research contribution.
\end{description}

The pipeline sweeps across 63 hyperparameter configurations (3
label-smoothing values $\times$ 3 learning rates $\times$ 7 random seeds),
trains three backbone architectures (ConvNeXt-Tiny, DenseNet-121,
InceptionV3) under each configuration, and automatically evaluates each
result using accuracy metrics and a full suite of explainability methods
(GradCAM, LIME, SHAP, DiCE, t-SNE).

A secondary output is a demo notebook (\texttt{Stacked\_Model\_Colab.ipynb})
that loads the best-performing pre-trained models and runs ensemble
inference on a single bug-bite image.

\textbf{Important:} This system is a research prototype. Its outputs are
not suitable for clinical or diagnostic use.

% ============================================================
\subsection{Use Scenario 1: Running the Automated Training Pipeline}
\label{sec:ug-pipeline}
% ============================================================

This scenario describes how to run the full hyperparameter sweep from start
to results. A dedicated GPU machine is required.

\subsubsection{Prerequisites}

\begin{itemize}
    \item A CUDA-capable GPU (8\,GB VRAM minimum, 16\,GB or more
          recommended).
    \item Python 3.x with the required packages installed:
\begin{lstlisting}[language=bash]
pip install torch torchvision timm numpy Pillow opencv-python \
            scikit-learn matplotlib tqdm shap lime scikit-image dice-ml
\end{lstlisting}
    \item The project repository cloned locally.
    \item The bug-bite dataset (included in the repository under
          \texttt{Bug-Data/Multiclass\_Bug\_data/}).
    \item The cyclone dataset (fetched via the provided script, or
          available from \texttt{ishay.yulzary@gmail.com}).
\end{itemize}

\subsubsection{Step 1: Fetch the Cyclone Dataset}

From the repository root, run:

\begin{lstlisting}[language=bash]
# Full dataset: 4000 images per category (~4 hours)
python miscellaneous_code/fetch_cyclone_dataset.py

# Quick trial: 5 images per category (for testing the pipeline)
python miscellaneous_code/fetch_cyclone_dataset.py --trial
\end{lstlisting}

The download is resumable: if interrupted, re-run the same command and it
continues from where it stopped.

\subsubsection{Step 2: Configure Dataset Paths}

Open \texttt{miscellaneous\_code/run\_experiments.py} and set the four
path constants at the top of the file to point to your local dataset
locations:

\begin{lstlisting}[language=Python]
BUG_TRAIN = '/path/to/bug_data/train'
BUG_VAL   = '/path/to/bug_data/val'
TC_TRAIN  = '/path/to/cyclone_data_split/train'
TC_VAL    = '/path/to/cyclone_data_split/val'
\end{lstlisting}

\subsubsection{Step 3: Run the Sweep}

\begin{lstlisting}[language=bash]
python miscellaneous_code/run_experiments.py
\end{lstlisting}

The sweep iterates over all 63 configurations. For each configuration it
runs cyclone pre-training, bug-bite fine-tuning, GradCAM evaluation, and
the full XAI suite in sequence. Check progress at any time without
interrupting:

\begin{lstlisting}[language=bash]
python miscellaneous_code/run_experiments.py --status
\end{lstlisting}

The sweep can be paused safely at any point by pressing \textbf{Ctrl-C}
or creating the file \texttt{results/PAUSE}. Re-running the same command
resumes from where it stopped. Interrupted configurations are automatically
recovered on the next run.

\subsubsection{Step 4: Review Results}

After the sweep completes (or at any point during it), results are
available in the following locations:

\begin{description}
    \item[\texttt{results/experiment\_results.json}]
          Per-model, per-configuration validation accuracy, macro F1,
          precision, and recall for all completed configurations.
    \item[\texttt{results/xai/\{run-id\}/ensemble\_metrics.json}]
          Ensemble accuracy for each configuration (soft-vote across all
          three architectures, Control and TC compared directly).
    \item[\texttt{results/xai/\{run-id\}/silhouette.json}]
          Silhouette scores on backbone embeddings, measuring class
          separability in the learned feature space.
    \item[\texttt{results/plots/\{run-id\}/}]
          Training curves, per-class F1, and ROC curves.
    \item[\texttt{results/xai/\{run-id\}/}]
          t-SNE projections, GradCAM overlays, LIME visualisations, and
          SHAP attribution maps.
    \item[\texttt{results/checkpoints/}]
          Saved model weights for all completed configurations.
\end{description}

\textbf{Summary of published results:} TC pre-training matches or exceeds
the Control baseline across all tested configurations. The best TC ensemble
accuracy is \textbf{83.27\%} (mean over 7 seeds, ls=0.10,
lr=$5\times10^{-6}$) vs.\ a Control mean of \textbf{81.83\%}
($+1.45$ percentage points). The best single-run TC accuracy is
\textbf{86.08\%}.

% ============================================================
\subsection{Use Scenario 2: Inference via the Demo Notebook}
\label{sec:ug-inference}
% ============================================================

This scenario describes how to classify a single bug-bite image using the
pre-trained ensemble on Google Colab. No local installation or GPU is
required.

\subsubsection{Prerequisites}

\begin{enumerate}
    \item A Google account.
    \item The file \texttt{project\_book\_results.tar} placed in your
          Google Drive. Contact \texttt{ishay.yulzary@gmail.com} to
          request access. Recommended location:
\begin{lstlisting}
MyDrive/Capstone/project_book_results.tar
\end{lstlisting}
    \item A bug-bite image in JPEG or PNG format.
\end{enumerate}

\subsubsection{Step 1: Open the Notebook}

\begin{enumerate}
    \item Go to \url{https://colab.research.google.com}.
    \item Select \textbf{File} $\rightarrow$ \textbf{Open notebook}
          $\rightarrow$ \textbf{GitHub}.
    \item Paste the project repository URL and select
          \texttt{Stacked\_Model\_Colab.ipynb}.
\end{enumerate}

\subsubsection{Step 2: Install Dependencies}

Run \textbf{Section 1}. Installs all required packages. Takes approximately
one minute and only runs once per Colab session.

\subsubsection{Step 3: Mount Drive and Verify Setup}

Run \textbf{Section 2}. Authorise Colab to access your Google Drive.
Confirm the output reads:

\begin{lstlisting}
PROJECT_ROOT exists: True
Archive exists:      True
\end{lstlisting}

If your archive is stored at a different Drive path, update
\texttt{PROJECT\_ROOT} directly inside the Section 2 cell before running:

\begin{lstlisting}[language=Python]
PROJECT_ROOT = '/content/drive/MyDrive/YourFolder'
\end{lstlisting}

This must be changed inside the cell itself, not in a separate cell.

\subsubsection{Step 4: Extract Weights and Load Models}

Run \textbf{Section 3} to extract the archive (2--5 minutes on first run;
skipped automatically on subsequent runs). Then run \textbf{Section 4}
and \textbf{Section 5} to load the six pre-trained models. Successful
loading prints:

\begin{lstlisting}
All models loaded.
\end{lstlisting}

\subsubsection{Step 5: Run Inference}

In \textbf{Section 7}, set \texttt{IMG\_PATH} to your image:

\textbf{Option A: Image on Google Drive.}
\begin{lstlisting}[language=Python]
IMG_PATH = '/content/drive/MyDrive/your_image.jpg'
\end{lstlisting}

\textbf{Option B: Upload directly to Colab.}
\begin{lstlisting}[language=Python]
from google.colab import files
up = files.upload()
IMG_PATH = list(up.keys())[0]
\end{lstlisting}

Run the cell. The output shows the predicted bug-bite category and
a probability distribution over all five classes for both the Control
and TC ensembles side by side.

\subsubsection{Step 6: Generate Visual Explanations}

\textbf{GradCAM (Section 8):} produces a $3 \times 2$ grid of attention
heatmaps (one row per architecture, one column per pipeline). Warm-coloured
regions had the strongest influence on the prediction.

\textbf{LIME (Section 9):} produces superpixel attribution maps. Green
regions support the predicted class; red regions contradict it. Takes
2--5 minutes on a CPU runtime.

% ============================================================
\subsection{Contact}
\label{sec:ug-contact}
% ============================================================

For access to model weights, datasets, or the environment specification:

\medskip
\noindent\textbf{Ishay Yulzary} --- \texttt{ishay.yulzary@gmail.com}

% --- End of appendix (remove if standalone) ---
\end{document}

%
%  Or compile standalone:
%    Uncomment the preamble block below, and add \end{document} at the end.
% ============================================================

% --- Standalone preamble (remove if including in main document) ----------
\documentclass[12pt,a4paper]{article}
\usepackage[utf8]{inputenc}
\usepackage[T1]{fontenc}
\usepackage{geometry}
\geometry{margin=2.5cm}
\usepackage{hyperref}
\usepackage{listings}
\usepackage{xcolor}
\usepackage{booktabs}
\usepackage{array}
\usepackage{enumitem}
\usepackage{parskip}
\lstset{basicstyle=\ttfamily\small, breaklines=true, frame=single,
        backgroundcolor=\color{gray!8}}
\begin{document}
\title{User Guide\\
       \large Bug-Bite Classification via Tropical Cyclone Transfer Learning\\
       Braude College of Engineering --- Software Engineering Capstone}
\author{Ishay Yulzary, Nathan Trostianitser}
\date{}
\maketitle
\tableofcontents
\newpage
% -------------------------------------------------------------------------

\section{User Guide}
\label{app:user-guide}

This guide covers two use scenarios for the system:

\begin{enumerate}
    \item \textbf{Running the automated training pipeline:} execute the
          full hyperparameter sweep to train and evaluate all model
          configurations. This is the primary use of the system.
    \item \textbf{Running inference via the demo notebook:} classify a
          single bug-bite image using the pre-trained ensemble on Google
          Colab, without any local setup.
\end{enumerate}

% ============================================================
\subsection{What the System Does}
\label{sec:ug-overview}
% ============================================================

The primary output of this project is an automated machine learning
pipeline that trains and evaluates an ensemble of neural network classifiers
under two paradigms:

\begin{description}
    \item[Control] Models trained directly from ImageNet weights on the
          bug-bite dataset (single-stage transfer).
    \item[TC (Transfer Classification)] Models first pre-trained on tropical
          cyclone satellite imagery, then fine-tuned on the bug-bite dataset
          (two-stage transfer). This is the primary research contribution.
\end{description}

The pipeline sweeps across 63 hyperparameter configurations (3
label-smoothing values $\times$ 3 learning rates $\times$ 7 random seeds),
trains three backbone architectures (ConvNeXt-Tiny, DenseNet-121,
InceptionV3) under each configuration, and automatically evaluates each
result using accuracy metrics and a full suite of explainability methods
(GradCAM, LIME, SHAP, DiCE, t-SNE).

A secondary output is a demo notebook (\texttt{Stacked\_Model\_Colab.ipynb})
that loads the best-performing pre-trained models and runs ensemble
inference on a single bug-bite image.

\textbf{Important:} This system is a research prototype. Its outputs are
not suitable for clinical or diagnostic use.

% ============================================================
\subsection{Use Scenario 1: Running the Automated Training Pipeline}
\label{sec:ug-pipeline}
% ============================================================

This scenario describes how to run the full hyperparameter sweep from start
to results. A dedicated GPU machine is required.

\subsubsection{Prerequisites}

\begin{itemize}
    \item A CUDA-capable GPU (8\,GB VRAM minimum, 16\,GB or more
          recommended).
    \item Python 3.x with the required packages installed:
\begin{lstlisting}[language=bash]
pip install torch torchvision timm numpy Pillow opencv-python \
            scikit-learn matplotlib tqdm shap lime scikit-image dice-ml
\end{lstlisting}
    \item The project repository cloned locally.
    \item The bug-bite dataset (included in the repository under
          \texttt{Bug-Data/Multiclass\_Bug\_data/}).
    \item The cyclone dataset (fetched via the provided script, or
          available from \texttt{ishay.yulzary@gmail.com}).
\end{itemize}

\subsubsection{Step 1: Fetch the Cyclone Dataset}

From the repository root, run:

\begin{lstlisting}[language=bash]
# Full dataset: 4000 images per category (~4 hours)
python miscellaneous_code/fetch_cyclone_dataset.py

# Quick trial: 5 images per category (for testing the pipeline)
python miscellaneous_code/fetch_cyclone_dataset.py --trial
\end{lstlisting}

The download is resumable: if interrupted, re-run the same command and it
continues from where it stopped.

\subsubsection{Step 2: Configure Dataset Paths}

Open \texttt{miscellaneous\_code/run\_experiments.py} and set the four
path constants at the top of the file to point to your local dataset
locations:

\begin{lstlisting}[language=Python]
BUG_TRAIN = '/path/to/bug_data/train'
BUG_VAL   = '/path/to/bug_data/val'
TC_TRAIN  = '/path/to/cyclone_data_split/train'
TC_VAL    = '/path/to/cyclone_data_split/val'
\end{lstlisting}

\subsubsection{Step 3: Run the Sweep}

\begin{lstlisting}[language=bash]
python miscellaneous_code/run_experiments.py
\end{lstlisting}

The sweep iterates over all 63 configurations. For each configuration it
runs cyclone pre-training, bug-bite fine-tuning, GradCAM evaluation, and
the full XAI suite in sequence. Check progress at any time without
interrupting:

\begin{lstlisting}[language=bash]
python miscellaneous_code/run_experiments.py --status
\end{lstlisting}

The sweep can be paused safely at any point by pressing \textbf{Ctrl-C}
or creating the file \texttt{results/PAUSE}. Re-running the same command
resumes from where it stopped. Interrupted configurations are automatically
recovered on the next run.

\subsubsection{Step 4: Review Results}

After the sweep completes (or at any point during it), results are
available in the following locations:

\begin{description}
    \item[\texttt{results/experiment\_results.json}]
          Per-model, per-configuration validation accuracy, macro F1,
          precision, and recall for all completed configurations.
    \item[\texttt{results/xai/\{run-id\}/ensemble\_metrics.json}]
          Ensemble accuracy for each configuration (soft-vote across all
          three architectures, Control and TC compared directly).
    \item[\texttt{results/xai/\{run-id\}/silhouette.json}]
          Silhouette scores on backbone embeddings, measuring class
          separability in the learned feature space.
    \item[\texttt{results/plots/\{run-id\}/}]
          Training curves, per-class F1, and ROC curves.
    \item[\texttt{results/xai/\{run-id\}/}]
          t-SNE projections, GradCAM overlays, LIME visualisations, and
          SHAP attribution maps.
    \item[\texttt{results/checkpoints/}]
          Saved model weights for all completed configurations.
\end{description}

\textbf{Summary of published results:} TC pre-training matches or exceeds
the Control baseline across all tested configurations. The best TC ensemble
accuracy is \textbf{83.27\%} (mean over 7 seeds, ls=0.10,
lr=$5\times10^{-6}$) vs.\ a Control mean of \textbf{81.83\%}
($+1.45$ percentage points). The best single-run TC accuracy is
\textbf{86.08\%}.

% ============================================================
\subsection{Use Scenario 2: Inference via the Demo Notebook}
\label{sec:ug-inference}
% ============================================================

This scenario describes how to classify a single bug-bite image using the
pre-trained ensemble on Google Colab. No local installation or GPU is
required.

\subsubsection{Prerequisites}

\begin{enumerate}
    \item A Google account.
    \item The file \texttt{project\_book\_results.tar} placed in your
          Google Drive. Contact \texttt{ishay.yulzary@gmail.com} to
          request access. Recommended location:
\begin{lstlisting}
MyDrive/Capstone/project_book_results.tar
\end{lstlisting}
    \item A bug-bite image in JPEG or PNG format.
\end{enumerate}

\subsubsection{Step 1: Open the Notebook}

\begin{enumerate}
    \item Go to \url{https://colab.research.google.com}.
    \item Select \textbf{File} $\rightarrow$ \textbf{Open notebook}
          $\rightarrow$ \textbf{GitHub}.
    \item Paste the project repository URL and select
          \texttt{Stacked\_Model\_Colab.ipynb}.
\end{enumerate}

\subsubsection{Step 2: Install Dependencies}

Run \textbf{Section 1}. Installs all required packages. Takes approximately
one minute and only runs once per Colab session.

\subsubsection{Step 3: Mount Drive and Verify Setup}

Run \textbf{Section 2}. Authorise Colab to access your Google Drive.
Confirm the output reads:

\begin{lstlisting}
PROJECT_ROOT exists: True
Archive exists:      True
\end{lstlisting}

If your archive is stored at a different Drive path, update
\texttt{PROJECT\_ROOT} directly inside the Section 2 cell before running:

\begin{lstlisting}[language=Python]
PROJECT_ROOT = '/content/drive/MyDrive/YourFolder'
\end{lstlisting}

This must be changed inside the cell itself, not in a separate cell.

\subsubsection{Step 4: Extract Weights and Load Models}

Run \textbf{Section 3} to extract the archive (2--5 minutes on first run;
skipped automatically on subsequent runs). Then run \textbf{Section 4}
and \textbf{Section 5} to load the six pre-trained models. Successful
loading prints:

\begin{lstlisting}
All models loaded.
\end{lstlisting}

\subsubsection{Step 5: Run Inference}

In \textbf{Section 7}, set \texttt{IMG\_PATH} to your image:

\textbf{Option A: Image on Google Drive.}
\begin{lstlisting}[language=Python]
IMG_PATH = '/content/drive/MyDrive/your_image.jpg'
\end{lstlisting}

\textbf{Option B: Upload directly to Colab.}
\begin{lstlisting}[language=Python]
from google.colab import files
up = files.upload()
IMG_PATH = list(up.keys())[0]
\end{lstlisting}

Run the cell. The output shows the predicted bug-bite category and
a probability distribution over all five classes for both the Control
and TC ensembles side by side.

\subsubsection{Step 6: Generate Visual Explanations}

\textbf{GradCAM (Section 8):} produces a $3 \times 2$ grid of attention
heatmaps (one row per architecture, one column per pipeline). Warm-coloured
regions had the strongest influence on the prediction.

\textbf{LIME (Section 9):} produces superpixel attribution maps. Green
regions support the predicted class; red regions contradict it. Takes
2--5 minutes on a CPU runtime.

% ============================================================
\subsection{Contact}
\label{sec:ug-contact}
% ============================================================

For access to model weights, datasets, or the environment specification:

\medskip
\noindent\textbf{Ishay Yulzary} --- \texttt{ishay.yulzary@gmail.com}

% --- End of appendix (remove if standalone) ---
\end{document}

%
%  Or compile standalone:
%    Uncomment the preamble block below, and add \end{document} at the end.
% ============================================================

% --- Standalone preamble (remove if including in main document) ----------
\documentclass[12pt,a4paper]{article}
\usepackage[utf8]{inputenc}
\usepackage[T1]{fontenc}
\usepackage{geometry}
\geometry{margin=2.5cm}
\usepackage{hyperref}
\usepackage{listings}
\usepackage{xcolor}
\usepackage{booktabs}
\usepackage{array}
\usepackage{enumitem}
\usepackage{parskip}
\lstset{basicstyle=\ttfamily\small, breaklines=true, frame=single,
        backgroundcolor=\color{gray!8}}
\begin{document}
\title{User Guide\\
       \large Bug-Bite Classification via Tropical Cyclone Transfer Learning\\
       Braude College of Engineering --- Software Engineering Capstone}
\author{Ishay Yulzary, Nathan Trostianitser}
\date{}
\maketitle
\tableofcontents
\newpage
% -------------------------------------------------------------------------

\section{User Guide}
\label{app:user-guide}

This guide covers two use scenarios for the system:

\begin{enumerate}
    \item \textbf{Running the automated training pipeline:} execute the
          full hyperparameter sweep to train and evaluate all model
          configurations. This is the primary use of the system.
    \item \textbf{Running inference via the demo notebook:} classify a
          single bug-bite image using the pre-trained ensemble on Google
          Colab, without any local setup.
\end{enumerate}

% ============================================================
\subsection{What the System Does}
\label{sec:ug-overview}
% ============================================================

The primary output of this project is an automated machine learning
pipeline that trains and evaluates an ensemble of neural network classifiers
under two paradigms:

\begin{description}
    \item[Control] Models trained directly from ImageNet weights on the
          bug-bite dataset (single-stage transfer).
    \item[TC (Transfer Classification)] Models first pre-trained on tropical
          cyclone satellite imagery, then fine-tuned on the bug-bite dataset
          (two-stage transfer). This is the primary research contribution.
\end{description}

The pipeline sweeps across 63 hyperparameter configurations (3
label-smoothing values $\times$ 3 learning rates $\times$ 7 random seeds),
trains three backbone architectures (ConvNeXt-Tiny, DenseNet-121,
InceptionV3) under each configuration, and automatically evaluates each
result using accuracy metrics and a full suite of explainability methods
(GradCAM, LIME, SHAP, DiCE, t-SNE).

A secondary output is a demo notebook (\texttt{Stacked\_Model\_Colab.ipynb})
that loads the best-performing pre-trained models and runs ensemble
inference on a single bug-bite image.

\textbf{Important:} This system is a research prototype. Its outputs are
not suitable for clinical or diagnostic use.

% ============================================================
\subsection{Use Scenario 1: Running the Automated Training Pipeline}
\label{sec:ug-pipeline}
% ============================================================

This scenario describes how to run the full hyperparameter sweep from start
to results. A dedicated GPU machine is required.

\subsubsection{Prerequisites}

\begin{itemize}
    \item A CUDA-capable GPU (8\,GB VRAM minimum, 16\,GB or more
          recommended).
    \item Python 3.x with the required packages installed:
\begin{lstlisting}[language=bash]
pip install torch torchvision timm numpy Pillow opencv-python \
            scikit-learn matplotlib tqdm shap lime scikit-image dice-ml
\end{lstlisting}
    \item The project repository cloned locally.
    \item The bug-bite dataset (included in the repository under
          \texttt{Bug-Data/Multiclass\_Bug\_data/}).
    \item The cyclone dataset (fetched via the provided script, or
          available from \texttt{ishay.yulzary@gmail.com}).
\end{itemize}

\subsubsection{Step 1: Fetch the Cyclone Dataset}

From the repository root, run:

\begin{lstlisting}[language=bash]
# Full dataset: 4000 images per category (~4 hours)
python miscellaneous_code/fetch_cyclone_dataset.py

# Quick trial: 5 images per category (for testing the pipeline)
python miscellaneous_code/fetch_cyclone_dataset.py --trial
\end{lstlisting}

The download is resumable: if interrupted, re-run the same command and it
continues from where it stopped.

\subsubsection{Step 2: Configure Dataset Paths}

Open \texttt{miscellaneous\_code/run\_experiments.py} and set the four
path constants at the top of the file to point to your local dataset
locations:

\begin{lstlisting}[language=Python]
BUG_TRAIN = '/path/to/bug_data/train'
BUG_VAL   = '/path/to/bug_data/val'
TC_TRAIN  = '/path/to/cyclone_data_split/train'
TC_VAL    = '/path/to/cyclone_data_split/val'
\end{lstlisting}

\subsubsection{Step 3: Run the Sweep}

\begin{lstlisting}[language=bash]
python miscellaneous_code/run_experiments.py
\end{lstlisting}

The sweep iterates over all 63 configurations. For each configuration it
runs cyclone pre-training, bug-bite fine-tuning, GradCAM evaluation, and
the full XAI suite in sequence. Check progress at any time without
interrupting:

\begin{lstlisting}[language=bash]
python miscellaneous_code/run_experiments.py --status
\end{lstlisting}

The sweep can be paused safely at any point by pressing \textbf{Ctrl-C}
or creating the file \texttt{results/PAUSE}. Re-running the same command
resumes from where it stopped. Interrupted configurations are automatically
recovered on the next run.

\subsubsection{Step 4: Review Results}

After the sweep completes (or at any point during it), results are
available in the following locations:

\begin{description}
    \item[\texttt{results/experiment\_results.json}]
          Per-model, per-configuration validation accuracy, macro F1,
          precision, and recall for all completed configurations.
    \item[\texttt{results/xai/\{run-id\}/ensemble\_metrics.json}]
          Ensemble accuracy for each configuration (soft-vote across all
          three architectures, Control and TC compared directly).
    \item[\texttt{results/xai/\{run-id\}/silhouette.json}]
          Silhouette scores on backbone embeddings, measuring class
          separability in the learned feature space.
    \item[\texttt{results/plots/\{run-id\}/}]
          Training curves, per-class F1, and ROC curves.
    \item[\texttt{results/xai/\{run-id\}/}]
          t-SNE projections, GradCAM overlays, LIME visualisations, and
          SHAP attribution maps.
    \item[\texttt{results/checkpoints/}]
          Saved model weights for all completed configurations.
\end{description}

\textbf{Summary of published results:} TC pre-training matches or exceeds
the Control baseline across all tested configurations. The best TC ensemble
accuracy is \textbf{83.27\%} (mean over 7 seeds, ls=0.10,
lr=$5\times10^{-6}$) vs.\ a Control mean of \textbf{81.83\%}
($+1.45$ percentage points). The best single-run TC accuracy is
\textbf{86.08\%}.

% ============================================================
\subsection{Use Scenario 2: Inference via the Demo Notebook}
\label{sec:ug-inference}
% ============================================================

This scenario describes how to classify a single bug-bite image using the
pre-trained ensemble on Google Colab. No local installation or GPU is
required.

\subsubsection{Prerequisites}

\begin{enumerate}
    \item A Google account.
    \item The file \texttt{project\_book\_results.tar} placed in your
          Google Drive. Contact \texttt{ishay.yulzary@gmail.com} to
          request access. Recommended location:
\begin{lstlisting}
MyDrive/Capstone/project_book_results.tar
\end{lstlisting}
    \item A bug-bite image in JPEG or PNG format.
\end{enumerate}

\subsubsection{Step 1: Open the Notebook}

\begin{enumerate}
    \item Go to \url{https://colab.research.google.com}.
    \item Select \textbf{File} $\rightarrow$ \textbf{Open notebook}
          $\rightarrow$ \textbf{GitHub}.
    \item Paste the project repository URL and select
          \texttt{Stacked\_Model\_Colab.ipynb}.
\end{enumerate}

\subsubsection{Step 2: Install Dependencies}

Run \textbf{Section 1}. Installs all required packages. Takes approximately
one minute and only runs once per Colab session.

\subsubsection{Step 3: Mount Drive and Verify Setup}

Run \textbf{Section 2}. Authorise Colab to access your Google Drive.
Confirm the output reads:

\begin{lstlisting}
PROJECT_ROOT exists: True
Archive exists:      True
\end{lstlisting}

If your archive is stored at a different Drive path, update
\texttt{PROJECT\_ROOT} directly inside the Section 2 cell before running:

\begin{lstlisting}[language=Python]
PROJECT_ROOT = '/content/drive/MyDrive/YourFolder'
\end{lstlisting}

This must be changed inside the cell itself, not in a separate cell.

\subsubsection{Step 4: Extract Weights and Load Models}

Run \textbf{Section 3} to extract the archive (2--5 minutes on first run;
skipped automatically on subsequent runs). Then run \textbf{Section 4}
and \textbf{Section 5} to load the six pre-trained models. Successful
loading prints:

\begin{lstlisting}
All models loaded.
\end{lstlisting}

\subsubsection{Step 5: Run Inference}

In \textbf{Section 7}, set \texttt{IMG\_PATH} to your image:

\textbf{Option A: Image on Google Drive.}
\begin{lstlisting}[language=Python]
IMG_PATH = '/content/drive/MyDrive/your_image.jpg'
\end{lstlisting}

\textbf{Option B: Upload directly to Colab.}
\begin{lstlisting}[language=Python]
from google.colab import files
up = files.upload()
IMG_PATH = list(up.keys())[0]
\end{lstlisting}

Run the cell. The output shows the predicted bug-bite category and
a probability distribution over all five classes for both the Control
and TC ensembles side by side.

\subsubsection{Step 6: Generate Visual Explanations}

\textbf{GradCAM (Section 8):} produces a $3 \times 2$ grid of attention
heatmaps (one row per architecture, one column per pipeline). Warm-coloured
regions had the strongest influence on the prediction.

\textbf{LIME (Section 9):} produces superpixel attribution maps. Green
regions support the predicted class; red regions contradict it. Takes
2--5 minutes on a CPU runtime.

% ============================================================
\subsection{Contact}
\label{sec:ug-contact}
% ============================================================

For access to model weights, datasets, or the environment specification:

\medskip
\noindent\textbf{Ishay Yulzary} --- \texttt{ishay.yulzary@gmail.com}

% --- End of appendix (remove if standalone) ---
\end{document}

%
%  Or compile standalone:
%    Uncomment the preamble block below, and add \end{document} at the end.
% ============================================================

% --- Standalone preamble (remove if including in main document) ----------
\documentclass[12pt,a4paper]{article}
\usepackage[utf8]{inputenc}
\usepackage[T1]{fontenc}
\usepackage{geometry}
\geometry{margin=2.5cm}
\usepackage{hyperref}
\usepackage{listings}
\usepackage{xcolor}
\usepackage{booktabs}
\usepackage{array}
\usepackage{enumitem}
\usepackage{parskip}
\lstset{basicstyle=\ttfamily\small, breaklines=true, frame=single,
        backgroundcolor=\color{gray!8}}
\begin{document}
\title{User Guide\\
       \large Bug-Bite Classification via Tropical Cyclone Transfer Learning\\
       Braude College of Engineering --- Software Engineering Capstone}
\author{Ishay Yulzary, Nathan Trostianitser}
\date{}
\maketitle
\tableofcontents
\newpage
% -------------------------------------------------------------------------

\section{User Guide}
\label{app:user-guide}

This guide covers two use scenarios for the system:

\begin{enumerate}
    \item \textbf{Running the automated training pipeline:} execute the
          full hyperparameter sweep to train and evaluate all model
          configurations. This is the primary use of the system.
    \item \textbf{Running inference via the demo notebook:} classify a
          single bug-bite image using the pre-trained ensemble on Google
          Colab, without any local setup.
\end{enumerate}

% ============================================================
\subsection{What the System Does}
\label{sec:ug-overview}
% ============================================================

The primary output of this project is an automated machine learning
pipeline that trains and evaluates an ensemble of neural network classifiers
under two paradigms:

\begin{description}
    \item[Control] Models trained directly from ImageNet weights on the
          bug-bite dataset (single-stage transfer).
    \item[TC (Transfer Classification)] Models first pre-trained on tropical
          cyclone satellite imagery, then fine-tuned on the bug-bite dataset
          (two-stage transfer). This is the primary research contribution.
\end{description}

The pipeline sweeps across 63 hyperparameter configurations (3
label-smoothing values $\times$ 3 learning rates $\times$ 7 random seeds),
trains three backbone architectures (ConvNeXt-Tiny, DenseNet-121,
InceptionV3) under each configuration, and automatically evaluates each
result using accuracy metrics and a full suite of explainability methods
(GradCAM, LIME, SHAP, DiCE, t-SNE).

A secondary output is a demo notebook (\texttt{Stacked\_Model\_Colab.ipynb})
that loads the best-performing pre-trained models and runs ensemble
inference on a single bug-bite image.

\textbf{Important:} This system is a research prototype. Its outputs are
not suitable for clinical or diagnostic use.

% ============================================================
\subsection{Use Scenario 1: Running the Automated Training Pipeline}
\label{sec:ug-pipeline}
% ============================================================

This scenario describes how to run the full hyperparameter sweep from start
to results. A dedicated GPU machine is required.

\subsubsection{Prerequisites}

\begin{itemize}
    \item A CUDA-capable GPU (8\,GB VRAM minimum, 16\,GB or more
          recommended).
    \item Python 3.x with the required packages installed:
\begin{lstlisting}[language=bash]
pip install torch torchvision timm numpy Pillow opencv-python \
            scikit-learn matplotlib tqdm shap lime scikit-image dice-ml
\end{lstlisting}
    \item The project repository cloned locally.
    \item The bug-bite dataset (included in the repository under
          \texttt{Bug-Data/Multiclass\_Bug\_data/}).
    \item The cyclone dataset (fetched via the provided script, or
          available from \texttt{ishay.yulzary@gmail.com}).
\end{itemize}

\subsubsection{Step 1: Fetch the Cyclone Dataset}

From the repository root, run:

\begin{lstlisting}[language=bash]
# Full dataset: 4000 images per category (~4 hours)
python miscellaneous_code/fetch_cyclone_dataset.py

# Quick trial: 5 images per category (for testing the pipeline)
python miscellaneous_code/fetch_cyclone_dataset.py --trial
\end{lstlisting}

The download is resumable: if interrupted, re-run the same command and it
continues from where it stopped.

\subsubsection{Step 2: Configure Dataset Paths}

Open \texttt{miscellaneous\_code/run\_experiments.py} and set the four
path constants at the top of the file to point to your local dataset
locations:

\begin{lstlisting}[language=Python]
BUG_TRAIN = '/path/to/bug_data/train'
BUG_VAL   = '/path/to/bug_data/val'
TC_TRAIN  = '/path/to/cyclone_data_split/train'
TC_VAL    = '/path/to/cyclone_data_split/val'
\end{lstlisting}

\subsubsection{Step 3: Run the Sweep}

\begin{lstlisting}[language=bash]
python miscellaneous_code/run_experiments.py
\end{lstlisting}

The sweep iterates over all 63 configurations. For each configuration it
runs cyclone pre-training, bug-bite fine-tuning, GradCAM evaluation, and
the full XAI suite in sequence. Check progress at any time without
interrupting:

\begin{lstlisting}[language=bash]
python miscellaneous_code/run_experiments.py --status
\end{lstlisting}

The sweep can be paused safely at any point by pressing \textbf{Ctrl-C}
or creating the file \texttt{results/PAUSE}. Re-running the same command
resumes from where it stopped. Interrupted configurations are automatically
recovered on the next run.

\subsubsection{Step 4: Review Results}

After the sweep completes (or at any point during it), results are
available in the following locations:

\begin{description}
    \item[\texttt{results/experiment\_results.json}]
          Per-model, per-configuration validation accuracy, macro F1,
          precision, and recall for all completed configurations.
    \item[\texttt{results/xai/\{run-id\}/ensemble\_metrics.json}]
          Ensemble accuracy for each configuration (soft-vote across all
          three architectures, Control and TC compared directly).
    \item[\texttt{results/xai/\{run-id\}/silhouette.json}]
          Silhouette scores on backbone embeddings, measuring class
          separability in the learned feature space.
    \item[\texttt{results/plots/\{run-id\}/}]
          Training curves, per-class F1, and ROC curves.
    \item[\texttt{results/xai/\{run-id\}/}]
          t-SNE projections, GradCAM overlays, LIME visualisations, and
          SHAP attribution maps.
    \item[\texttt{results/checkpoints/}]
          Saved model weights for all completed configurations.
\end{description}

\textbf{Summary of published results:} TC pre-training matches or exceeds
the Control baseline across all tested configurations. The best TC ensemble
accuracy is \textbf{83.27\%} (mean over 7 seeds, ls=0.10,
lr=$5\times10^{-6}$) vs.\ a Control mean of \textbf{81.83\%}
($+1.45$ percentage points). The best single-run TC accuracy is
\textbf{86.08\%}.

% ============================================================
\subsection{Use Scenario 2: Inference via the Demo Notebook}
\label{sec:ug-inference}
% ============================================================

This scenario describes how to classify a single bug-bite image using the
pre-trained ensemble on Google Colab. No local installation or GPU is
required.

\subsubsection{Prerequisites}

\begin{enumerate}
    \item A Google account.
    \item The file \texttt{project\_book\_results.tar} placed in your
          Google Drive. Contact \texttt{ishay.yulzary@gmail.com} to
          request access. Recommended location:
\begin{lstlisting}
MyDrive/Capstone/project_book_results.tar
\end{lstlisting}
    \item A bug-bite image in JPEG or PNG format.
\end{enumerate}

\subsubsection{Step 1: Open the Notebook}

\begin{enumerate}
    \item Go to \url{https://colab.research.google.com}.
    \item Select \textbf{File} $\rightarrow$ \textbf{Open notebook}
          $\rightarrow$ \textbf{GitHub}.
    \item Paste the project repository URL and select
          \texttt{Stacked\_Model\_Colab.ipynb}.
\end{enumerate}

\subsubsection{Step 2: Install Dependencies}

Run \textbf{Section 1}. Installs all required packages. Takes approximately
one minute and only runs once per Colab session.

\subsubsection{Step 3: Mount Drive and Verify Setup}

Run \textbf{Section 2}. Authorise Colab to access your Google Drive.
Confirm the output reads:

\begin{lstlisting}
PROJECT_ROOT exists: True
Archive exists:      True
\end{lstlisting}

If your archive is stored at a different Drive path, update
\texttt{PROJECT\_ROOT} directly inside the Section 2 cell before running:

\begin{lstlisting}[language=Python]
PROJECT_ROOT = '/content/drive/MyDrive/YourFolder'
\end{lstlisting}

This must be changed inside the cell itself, not in a separate cell.

\subsubsection{Step 4: Extract Weights and Load Models}

Run \textbf{Section 3} to extract the archive (2--5 minutes on first run;
skipped automatically on subsequent runs). Then run \textbf{Section 4}
and \textbf{Section 5} to load the six pre-trained models. Successful
loading prints:

\begin{lstlisting}
All models loaded.
\end{lstlisting}

\subsubsection{Step 5: Run Inference}

In \textbf{Section 7}, set \texttt{IMG\_PATH} to your image:

\textbf{Option A: Image on Google Drive.}
\begin{lstlisting}[language=Python]
IMG_PATH = '/content/drive/MyDrive/your_image.jpg'
\end{lstlisting}

\textbf{Option B: Upload directly to Colab.}
\begin{lstlisting}[language=Python]
from google.colab import files
up = files.upload()
IMG_PATH = list(up.keys())[0]
\end{lstlisting}

Run the cell. The output shows the predicted bug-bite category and
a probability distribution over all five classes for both the Control
and TC ensembles side by side.

\subsubsection{Step 6: Generate Visual Explanations}

\textbf{GradCAM (Section 8):} produces a $3 \times 2$ grid of attention
heatmaps (one row per architecture, one column per pipeline). Warm-coloured
regions had the strongest influence on the prediction.

\textbf{LIME (Section 9):} produces superpixel attribution maps. Green
regions support the predicted class; red regions contradict it. Takes
2--5 minutes on a CPU runtime.

% ============================================================
\subsection{Contact}
\label{sec:ug-contact}
% ============================================================

For access to model weights, datasets, or the environment specification:

\medskip
\noindent\textbf{Ishay Yulzary} --- \texttt{ishay.yulzary@gmail.com}

% --- End of appendix (remove if standalone) ---
\end{document}
